\section{Implications and a Scoped-Badge Study}\label{sec:agenda}
\noindent\textbf{Reading a badge.} We translate the scope contract into a proposed display such as ``call found, non-blocking behaviour unassessed'' for the first sentence in Fig.~\ref{fig:example}. Our planned study tests its effect on sentence judgments, motivated by evidence that automation cues can encourage misuse~\cite{parasuraman1997automation,bucinca2021trust}. We distinguish three properties: coverage counts the targets receiving a result, correctness concerns whether that result is right, and sentence sufficiency concerns whether its evidence covers the whole assertion. Improving one need not improve the others.

A false sentence with a valid subclaim makes this distinction concrete. Consider the hypothetical ``This calls \texttt{read} and never waits,'' when the call exists but can wait. Highlighting the call supplies valid evidence for the subclaim, yet leaves the whole sentence unjustified. This case guides our planned study.

\noindent\textbf{Primary comparison.} We will compare scoped badges, which display the supported proposition and unassessed content, with generic ``checked'' badges on these false partial-support sentences. Our primary outcome is correct rejection of the whole sentence. A third condition presents prose with code links. All conditions provide identical code/evidence access.

Before recruitment, we require two independent humans to validate sentence truth, evidence validity and evidence sufficiency separately. Each participant will see 24 distinct items, eight per condition. We use balanced Latin-square assignment to rotate item blocks across interfaces without repeating an item within a participant. Each condition has four true and four false sentences. The false items comprise two valid partial-support, one misleading-evidence and one no-evidence item. Alongside evidence comprehension, we will record clicks, time, confidence accuracy (Brier score: squared error of a probability judgment) and judgments of nearby unbadged statements~\cite{pennycook2020impliedtruth}. These secondary measures include clicks, which alone do not establish comprehension. Our artifact details allocation, ethics, consent, debriefing and registration requirements.

% Balance the final main-text page, leaving the references on their own page.
\balance
\noindent\textbf{Design sensitivity.} Our planned logistic mixed model estimates correct-rejection probability with participant and item baseline variation (crossed random intercepts). Interface condition, evidence category and their interaction let the badge effect vary by evidence kind. Our primary test checks for either an increase or decrease within partial-support items at the 5\% significance level.

We fitted this model with the R package \texttt{lme4} to 60,000 simulated datasets: five participant counts, four target odds ratios and three levels of participant/item variation, with 1,000 repetitions per setting. For rejection probability $p$, odds are $p/(1-p)$. An odds ratio (OR) of 1.65 multiplies them by 1.65 with participant/item effects fixed. Power, the fraction of simulations detecting that effect, is \RevPowerSmall{} with 36 participants and \RevPowerLargest{} with 216. With no badge effect (OR 1), false-positive rates are \RevNullRange{}. Monte Carlo standard error, uncertainty from finite repetitions, is at most 0.016. We count fits that fail to settle on stable estimates (converge) as non-detections and report them separately from singular fits, where estimated participant or item variation collapses to zero.

Our simulations hold badge effects constant across participants and items and exclude learning and dropout. We will set the recruitment target after agreeing on a minimum meaningful effect. Participant outcomes remain pending. The next evaluation phase requires independent labels, validated reference slots and a fresh corpus.
